%Format: Latex
\documentclass{article}
\setlength{\textwidth}{17cm}
\setlength{\textheight}{9in}
\setlength{\topmargin}{-1.5cm}
\setlength{\oddsidemargin}{0in}
\setlength{\evensidemargin}{0in}

\usepackage{graphicx, listings,color}
\usepackage{fontspec}
\usepackage{tikz}
\usepackage{listings}
\usepackage{url}

\definecolor{gray}{rgb}{0.4,0.4,0.4}
\definecolor{darkblue}{rgb}{0.0,0.0,0.6}
\definecolor{cyan}{rgb}{0.0,0.6,0.6}

\def\cw#1{\texttt{#1}}

\lstset{
  basicstyle=\ttfamily,
  columns=fullflexible,
  showstringspaces=false,
  commentstyle=\color{gray}\upshape
}

\lstdefinelanguage{XML}
{
  basicstyle=\ttfamily\color{darkblue}\bfseries,
  morestring=[b]",
  morestring=[s]{>}{<},
  morecomment=[s]{<?}{?>},
  stringstyle=\color{black},
  identifierstyle=\color{darkblue},
  keywordstyle=\color{cyan},
  morekeywords={xmlns,version,type}% list your attributes here
}


\definecolor{javared}{rgb}{0.6,0,0} % for strings
\definecolor{javagreen}{rgb}{0.25,0.5,0.35} % comments
\definecolor{javapurple}{rgb}{0.5,0,0.35} % keywords
\definecolor{javadocblue}{rgb}{0.25,0.35,0.75} % javadoc

\lstset{language=Java,
basicstyle=\ttfamily,
keywordstyle=\color{javapurple}\bfseries,
stringstyle=\color{javared},
commentstyle=\color{javagreen},
morecomment=[s][\color{javadocblue}]{/**}{*/},
%numbers=left,
%numberstyle=\tiny\color{black},
%stepnumber=1,
numbersep=12pt,
tabsize=4,
showspaces=false,
showstringspaces=false}


\usepackage{hyperref}
\usepackage{float}
\floatstyle{boxed} 
\restylefloat{figure}


\begin{document}

\def\codefont{
\fontspec{Courier New}
  \fontsize{9pt}{11pt}\selectfont}
\definecolor{codebgcolor}{HTML}{EDEDED}
\newenvironment{code}
{\begin{center}
    \begin{tikzpicture}
      \node [fill=codebgcolor,rounded corners=5pt]
      \bgroup
      \bgroup\codefont
      \begin{tabular}{l}}
      {\end{tabular}
      \egroup
      \egroup;
    \end{tikzpicture}
  \end{center}}


\title{Android Crash Course - Notes}
\author{Typed by Adam Matan, \texttt{adam@matan.name} \\
\texttt{http://matan.name}}
\date{Technion, Israel, Summer 2012}

\maketitle

\begin{abstract}
This is a summary of the 
 lectures given in the by Dr. Amnon Dekel and Mr. Haim Michael in the Android
\end{abstract}

\tableofcontents

\section{UI}
\subsection{Layouts}
A layout is an object that defines the way UI elements are ordered.

\subsubsection{Relative layout}

The \texttt{RelativeLayout} layout defines the position of each element in relation to other elements and the layout itself. For example:

An \texttt{Activity} is preferred over a \texttt{Layout} because it is reusable. Suppose that we create an activity for logging in into the Technion user page. 

\begin{itemize}
\item The \texttt{gravity} property
\end{itemize}

\subsection{Menus and Dialogs}

Each Android device has a menu button. Older devices have a physical one, while newer devices have a touch area. The menu button is important because the user expects to interact with menu items like \texttt{exit} or  \texttt{settings}. Therefore, an application with a menu is often more intuitive. Each application is associated with a default menu, which has a \texttt{settings} item that does nothing.

A menu consists of menu items and submenus. The submenus are limited to one hierarchy level (i.e., there are no sub-sub menus). The items can be either associated with a \texttt{OnMenuItemClickListener}, or with an explicit \texttt{Intent}.

Menu items are not limited to text;  An icon can greatly increase their readability.

\begin{figure}[h!]
  \centering
      \includegraphics[width=0.5\textwidth]{img/icons.jpeg}
  \caption{Combining Icons and Text in a menu. This example is not taken from an Android environment.}
\end{figure}

\section{Activity Data}
\subsection{MIME-type}
\href{http://en.wikipedia.org/wiki/MIME}{Multipurpose Internet Mail Extensions (MIME)} is an Internet format that defines data type for files. It was originally developed for email and web resources. MIME is used to define the data that is stored in an Intent.

\begin{code} \begin{lstlisting}[language=Java]
helloLanguagesIntent.setType("vnd.abelski.hello.translate/vnd.dictionary");
\end{lstlisting} \end{code}

\section{Sensors}

\subsection{Basics}

getSystemService() is a method that "Return the handle to a system-level service by name", like VIBRATOR\_SERVICE or SENSOR\_SERVICE.

\begin{code} \begin{lstlisting}[language=Java]
manager = (SensorManager) getSystemService(Context.SENSOR_SERVICE);
sensor = manager.getDefaultSensor(Sensor.TYPE_LIGHT);
\end{lstlisting} \end{code}


\subsection{Threads}
Threads in Android work pretty much the same as they work in Java.

\subsubsection{The Main (UI) Thread}
The UI Thread (AKA main thread) is a special thread that has exclusive control on UI elements. Blocking or delaying the main thread is a hideous idea, which will probably cause an "Application Not Responsive" error message. Therefore, Networking activity should always be done in a different thread. In the recent versions of Android, this behaviour is forced - networking can not be done from the main thread.

\subsubsection{Updating UI elements from other Threads}

An ordinary thread may need to update a UI element. This can happen, for example, if some search result returned from a network connection, and we would like to update a text box. This is done by sending a Runnable() object to the main thread, using either post() or runOnUiThread(). The main thread has a queue of Runnables, and it executes the run() methods.

\begin{code} \begin{lstlisting}[language=Java]
public final void runOnUiThread (Runnable action)
public boolean post(Runnable action)
\end{lstlisting} \end{code}

Alternatively, the AsyncTask class can be used to run code from the main thread in a non-blocking fashion. For example, if  we define DownloadFilesTask that extends AsyncTask, we can invoke DownloadFilesTask from the main thread:

\begin{code} \begin{lstlisting}[language=Java]
new DownloadFilesTask().execute(url1, url2, url3);
\end{lstlisting} \end{code}

\subsubsection{The Handler Class}
The Handler Class is a nice inter-Thread communication mechanism. Handlers can be used to pass Runnables and Messages. 



\section{Java Web Services and DOM}
JSON, XML, and REST and DOM services were covered briefly.

\subsection{Web Services in Android}

The Android environment utilizes Apache's \cw{HttpClient} HttpClient, which has three important classes: \cw{HttpGet}, \cw{HttpPost} and \cw{HttpResponse}. 

There are some noteworthy points regarding \cw{HttpPost}. Let's examine the code from the lecture slides:

\begin{code} \begin{lstlisting}[language=Java]
...
HttpClient client = new DefaultHttpClient();
HttpPost request = new HttpPost("http://mysite.com/upload.php");
List<NameValuePair> parameters =
   new ArrayList<NameValuePair>();
parameters.add(
   new BasicNameValuePair("param_name", "param_value"));
UrlEncodedFormEntity encodedEntity =
   new UrlEncodedFormEntity(parameters);
request.setEntity(encodedEntity);
HttpResponse response = client.execute(request);
in = new BufferedReader(
   new InputStreamReader(response.getEntity().getContent()));
...
\end{lstlisting} \end{code}

\begin{itemize}
\item{The \cw{POST} parameters are set using a \cw{List} of \cw{NameValuePair}s. }
\item{The URL is encoded using \cw{UrlEncodedFormEntity}. The encoding sets spaces and other strange characters to unambiguous  \href{http://en.wikipedia.org/wiki/Percent-encoding}{Percent Encoding}. }
\end{itemize}

\section{Database}
\subsection{SQLite}

\begin{center}
\includegraphics[width=0.3\textwidth]{img/sqlite_logo_wikipedia.eps}
\end{center}

SQLite is an open-source, cross-platform relational database stored in a single file; it is the native database solution used in Android. It is simple to use, but lacks support in advanced features like Stored Procedures.

\subsection{Initializing a SQLite database}
 In order to use SQLite, the caller app creates a class that inherits \cw{SQLiteOpenHelper}, and needs to implement two methods. The first, \cw{onCreate(SQLiteDatabase database)} defines the behaviour when the app is first installed and the database does not yet exist. 

The seconds method, \cw{onUpgrade(SQLiteDatabase db, int oldVersion, int newVersion)}, is executed when a previous version of the app is found. Each app should present a DB version coupled with its app version. \cw{onUpgrade} converts old database schema to new database schema, given the old DB version and new DB version.

\begin{code} \begin{lstlisting}[language=Java]
public class DatabaseHelper extends SQLiteOpenHelper{}
@Override
public void onCreate(SQLiteDatabase database)
@Override
public void onUpgrade(SQLiteDatabase db, int oldVersion, int newVersion)
\end{lstlisting} \end{code}

The \cw{execSQL()} method is used for creating the database structure. It supports any old SQL command, except for \cw{SELECT}, \cw{INSERT}, \cw{UPDATE} and \cw{DELETE}.

\subsection{SQLite sans SQL?}
The aforementioned SQL operations are supported via corresponding delegator methods:

\begin{center}
    \begin{tabular}{ | l | l |}
    \hline
    SQL Command & Equivalent method \\ \hline
    \cw{INSERT} & BLAH \\ \hline
    \cw{UPDATE} & BLAH \\ \hline
    \cw{DELETE} & BLAH \\ \hline
    \cw{SELECT} & BLAH \\ \hline
    \end{tabular}
\end{center}

\subsection{The \_ID Column}
\includegraphics[width=0.3\textwidth]{img/warning_sign.eps}
Every table

\subsection{Thread concerns}


\end{document}
